% 08_finish.tex - final checks, plots and notes
\section{Check and plot the final run}

\step{Check the water balance}
\term{t4_audit_c.txt}
\begin{tocheck}
\begin{itemize}
\item The reservoir lines: \cmd{delivered / sent} is 1, and no seepage is left undelivered.
\item The surface-store line and the recharge line are 0.
\item The cumulative residual is \V{t4:c:auditRes}, and there are \V{t4:c:auditUnconv} unconverged steps.
\end{itemize}
\end{tocheck}

\step{Plot the results}
\term{t4_plot.txt}
\fig{t4_gw_budget.png}{\cmd{gw\_budget.png}: the budget of the coupled model. The model's own well and boundary
(\emph{model's own packages}), water-table ET, and the storage release of the drain-down that follows the steady first
day.}
\fig[0.72]{t4_gw_map.png}{\cmd{gw\_map.png}: water-table depth below the model top, on the model's own cells (converted
from feet to metres).}

\section{Notes}
Raven never changes your model. \cmd{out/mf6/} is a working copy, with the stress periods split into days, the taken-over
packages commented out and Raven's packages added. \cmd{\$WS/model/finished} holds the files as they are at the end of
this workshop. The manual of the coupling (chapter ``Coupling an existing MODFLOW~6 model'') describes restarts and later
start dates, DISV grids, the water mover, steady-state models and the safeguards. It is in the RavenHydroFramework-extra
repository (see \cmd{bin/README.txt}), with a larger example, \cmd{examples/Liard\_existing\_model}.
